% !TEX program = xelatex
\documentclass[11pt,a4paper]{article}

\usepackage[UTF8]{ctex}
\usepackage[margin=2.2cm]{geometry}
\usepackage{amsmath, amssymb, amsthm}
\usepackage{listings}
\usepackage{xcolor}
\usepackage{tikz}
\usepackage{booktabs}
\usepackage{multirow}
\usepackage{hyperref}
\usetikzlibrary{positioning, arrows.meta, shapes.geometric, fit, calc, decorations.pathreplacing}

\hypersetup{
  colorlinks=true,
  linkcolor=blue!60!black,
  urlcolor=teal!70!black,
}

\lstdefinelanguage{Rust}{
  morekeywords={fn,let,mut,pub,struct,impl,use,mod,self,Self,Option,Some,None,return,if,else,for,while,loop,match,break,continue,as,where,trait,type,enum,move,ref,const,static,unsafe,extern,crate,super,dyn,async,await,box,in},
  sensitive=true,
  morecomment=[l]{//},
  morecomment=[s]{/*}{*/},
  morestring=[b]",
}

\lstset{
  language=Rust,
  basicstyle=\ttfamily\footnotesize,
  keywordstyle=\color{blue!60!black}\bfseries,
  commentstyle=\color{green!40!black}\itshape,
  stringstyle=\color{red!60!black},
  numbers=left,
  numberstyle=\tiny\color{gray},
  numbersep=8pt,
  frame=single,
  framerule=0.4pt,
  rulecolor=\color{gray!40},
  breaklines=true,
  breakatwhitespace=true,
  showstringspaces=false,
  tabsize=4,
  columns=fullflexible,
  keepspaces=true,
}

\title{\textbf{halfedge\,: 面朝向一致性与可定向性模块设计文档} \\ \large \texttt{src/orientation.rs}}
\author{halfedge 库}
\date{2026-07-04}

\begin{document}
\maketitle
\tableofcontents

\section{模块定位}

\texttt{orientation.rs} 提供面朝向一致性检测与修复能力，回答以下三个问题：

\begin{enumerate}
  \item 网格中所有相邻面是否朝向一致（\texttt{are\_normals\_consistent}）？
  \item 网格是否可定向（\texttt{is\_orientable}）？
  \item 如何修复不一致的朝向（\texttt{fix\_orientations}）？
\end{enumerate}

\subsection{API 总览}

\begin{center}
\renewcommand{\arraystretch}{1.18}
\begin{tabular}{@{}lll@{}}
  \toprule
  \textbf{类别} & \textbf{API} & \textbf{语义} \\
  \midrule
  检测 & \texttt{are\_normals\_consistent(mesh)} & 是否所有相邻面朝向一致 \\
  检测 & \texttt{is\_orientable(mesh)}           & 是否可定向（无 Möbius 矛盾） \\
  修复 & \texttt{fix\_orientations(\&mut mesh)}   & 翻转不一致的面，返回翻转数 \\
  \bottomrule
\end{tabular}
\end{center}

\section{背景：朝向一致与可定向}

\subsection{朝向一致（consistently oriented）}

设半边 $h$ 与其 twin $h^\star$ 分别属于面 $f_1, f_2$：
\[
  h: A \to B,\ \mathrm{face}(h) = f_1; \qquad
  h^\star: B \to A,\ \mathrm{face}(h^\star) = f_2.
\]
\textbf{朝向一致}要求 $h$ 与 $h^\star$ \textbf{方向相反}，即
$\mathrm{vertex}(h) \ne \mathrm{vertex}(h^\star)$。

若两者方向相同（$\mathrm{vertex}(h) = \mathrm{vertex}(h^\star)$），
则两面朝向「同向」，构成局部不一致。

\subsection{可定向（orientable）}

\textbf{可定向}是更弱的性质：网格能否为每个面分配一个朝向，使所有相邻面
朝向一致。不可定向网格（如 Möbius 带、Klein 瓶）存在矛盾分配，
无法通过翻转修复。

\begin{center}
\renewcommand{\arraystretch}{1.18}
\begin{tabular}{@{}lll@{}}
  \toprule
  \textbf{性质} & \textbf{检查对象} & \textbf{失败原因} \\
  \midrule
  一致  & 当前 \texttt{mesh} 状态       & 存在同向 twin 对 \\
  可定向 & BFS 遍历中的方向分配          & 存在矛盾（Möbius 结构） \\
  \bottomrule
\end{tabular}
\end{center}

\section{are\_normals\_consistent 算法}

遍历所有半边，检查 twin 对的方向：

\begin{lstlisting}
pub fn are_normals_consistent(mesh: &MeshStorage) -> bool {
    for he in mesh.halfedge_ids() {
        let h = match mesh.get_halfedge(he) { Some(h) => h, None => continue };
        let twin = match h.twin { Some(t) => t, None => continue };
        let twin_data = match mesh.get_halfedge(twin) { Some(t) => t, None => continue };
        // he.vertex == twin.vertex 表示同向 -> 不一致
        if h.vertex == twin_data.vertex {
            return false;
        }
    }
    true
}
\end{lstlisting}

\textbf{复杂度}：$O(E)$（遍历所有半边），$O(1)$ 额外空间。
\textbf{提前终止}：发现第一处同向即返回 \texttt{false}。

\section{is\_orientable 算法}

\subsection{思路}

对每个面连通分量做 BFS，从基准面开始传播朝向约束。若某条 twin 边两端
方向相同，则不可定向。

\begin{center}
\resizebox{\textwidth}{!}{%
\begin{tikzpicture}[
  node distance=0.5cm,
  proc/.style={draw, rounded corners=2pt, minimum width=10cm, minimum height=0.7cm, align=center, font=\small},
  dec/.style={diamond, draw, aspect=2.4, fill=orange!12, inner sep=1pt, font=\small, align=center, minimum width=3.5cm},
  op/.style={-Stealth, thick},
  startend/.style={draw, rounded corners=10pt, minimum width=2.6cm, minimum height=0.7cm, font=\small, fill=gray!12}
]
  \node[startend] (s) {开始：遍历所有面};
  \node[dec, below=of s] (d1) {面 $f$ 已访问？};
  \node[proc, right=1.6cm of d1, fill=blue!8] (skip) {跳过};
  \node[proc, below=of d1, fill=blue!8] (p1) {$f \to$ 队列，标记访问};
  \node[dec, below=of p1] (d2) {队列非空？};
  \node[proc, below=of d2, fill=yellow!15] (p2) {出队 $f$；遍历 $f$ 的所有半边 $h$};
  \node[dec, below=of p2] (d3) {twin $h^\star$ 有邻接面 $f'$？};
  \node[dec, below=of d3] (d4) {$h.\mathrm{vertex} = h^\star.\mathrm{vertex}$？};
  \node[proc, right=1.6cm of d4, fill=red!15] (bad) {返回 \texttt{false}（不可定向）};
  \node[proc, below=of d4, fill=green!12] (ok) {标记 $f'$ 访问；入队 $f'$};
  \node[startend, right=1.6cm of d2, fill=green!12] (e) {返回 \texttt{true}};
  \draw[op] (s) -- (d1);
  \draw[op] (d1) -- node[above, font=\scriptsize] {是} (skip);
  \draw[op] (d1) -- node[right, font=\scriptsize] {否} (p1);
  \draw[op] (p1) -- (d2);
  \draw[op] (d2) -- node[above, font=\scriptsize] {否} (e);
  \draw[op] (d2) -- node[right, font=\scriptsize] {是} (p2);
  \draw[op] (p2) -- (d3);
  \draw[op] (d3) -- node[right, font=\scriptsize] {是} (d4);
  \draw[op] (d4) -- node[above, font=\scriptsize] {是} (bad);
  \draw[op] (d4) -- node[right, font=\scriptsize] {否} (ok);
  \draw[op, dashed] (ok.west) -| ($(p2.west)+(-0.6,0)$) -- (p2.west);
  \draw[op, dashed] (skip.north) |- (e.north);
\end{tikzpicture}
}
\end{center}

\subsection{与 are\_normals\_consistent 的区别}

\texttt{are\_normals\_consistent} 是\textbf{当前状态}检测，仅判断是否存在
同向 twin 对；\texttt{is\_orientable} 是\textbf{可分配性}检测，
即使当前状态不一致，也可能通过翻转某些面来达到一致（即可定向但当前不一致）。

\section{fix\_orientations 算法}

\subsection{整体流程}

对每个连通分量，BFS 遍历，遇到与基准方向不一致的邻接面就翻转：

\begin{center}
\resizebox{\textwidth}{!}{%
\begin{tikzpicture}[
  node distance=0.5cm,
  proc/.style={draw, rounded corners=2pt, minimum width=11cm, minimum height=0.7cm, align=center, font=\small},
  dec/.style={diamond, draw, aspect=2.4, fill=orange!12, inner sep=1pt, font=\small, align=center, minimum width=3.5cm},
  op/.style={-Stealth, thick},
  startend/.style={draw, rounded corners=10pt, minimum width=2.6cm, minimum height=0.7cm, font=\small, fill=gray!12}
]
  \node[startend] (s) {开始};
  \node[proc, below=of s, fill=blue!8] (p1) {对每个未访问的 \texttt{seed}：\texttt{visited $\gets \{$seed$\}$}；\texttt{queue $\gets$ [seed]}};
  \node[dec, below=of p1] (d1) {队列非空？};
  \node[proc, below=of d1, fill=yellow!15] (p2) {$f \gets$ queue.pop\_front()；遍历 $f$ 的所有半边 $h$};
  \node[dec, below=of p2] (d2) {邻接面 $f'$ 未访问且 $h.\mathrm{vertex} = h^\star.\mathrm{vertex}$？};
  \node[proc, right=1.6cm of d2, fill=red!12] (yes) {翻转 $f'$；\texttt{flipped $+{=}1$}};
  \node[proc, below=of d2, fill=green!12] (no) {标记 $f'$ 访问；入队 $f'$};
  \node[startend, right=1.6cm of d1, fill=green!12] (e) {返回 \texttt{flipped}};
  \draw[op] (s) -- (p1);
  \draw[op] (p1) -- (d1);
  \draw[op] (d1) -- node[right, font=\scriptsize] {是} (p2);
  \draw[op] (d1) -- node[above, font=\scriptsize] {否} (e);
  \draw[op] (p2) -- (d2);
  \draw[op] (d2) -- node[above, font=\scriptsize] {是} (yes);
  \draw[op] (d2) -- node[right, font=\scriptsize] {否} (no);
  \draw[op, dashed] (yes.east) -- ++(0.4,0) |- (no.east);
  \draw[op, dashed] (no.west) -| ($(p2.west)+(-0.6,0)$) -- (p2.west);
\end{tikzpicture}
}
\end{center}

\subsection{flip\_face\_orientation：翻转单个三角面}

对三角面 $[v_0, v_1, v_2]$，半边环 $h_0 \to h_1 \to h_2$，每条半边指向
对应顶点。翻转操作：

\begin{enumerate}
  \item \textbf{交换 vertex}：$h_0.\mathrm{vertex} \gets v_2$，
        $h_1.\mathrm{vertex} \gets v_0$，
        $h_2.\mathrm{vertex} \gets v_1$（每条半边指向原 next 的起点）；
  \item \textbf{交换 next/prev}：$h_i.\mathrm{next} \leftrightarrow h_i.\mathrm{prev}$。
\end{enumerate}

效果：环顺序由 $v_0 \to v_1 \to v_2$ 变为 $v_2 \to v_1 \to v_0$（CCW $\leftrightarrow$ CW）。

\subsection{限制}

\begin{itemize}
  \item 仅支持\textbf{三角面}（$k = 3$），\texttt{flip\_face\_orientation}
        内部对 $k \ne 3$ 仅交换 next/prev 但不调整 vertex，可能产生错误结果；
        本项目所有网格为三角网格，故无问题；
  \item 不可定向网格无法通过翻转达到一致，\texttt{fix\_orientations} 不报错，
        但剩余矛盾仍存在。调用方应在调用前用 \texttt{is\_orientable} 检查。
\end{itemize}

\section{复杂度}

\begin{center}
\renewcommand{\arraystretch}{1.18}
\begin{tabular}{@{}lll@{}}
  \toprule
  \textbf{API} & \textbf{时间复杂度} & \textbf{空间复杂度} \\
  \midrule
  \texttt{are\_normals\_consistent} & $O(E)$ & $O(1)$ \\
  \texttt{is\_orientable}           & $O(V + F)$ & $O(F)$（visited + 队列） \\
  \texttt{fix\_orientations}        & $O(V + F)$ & $O(F)$ \\
  \bottomrule
\end{tabular}
\end{center}

\section{退化守卫}

\begin{itemize}
  \item \texttt{are\_normals\_consistent}：\texttt{get\_halfedge} 返回 \texttt{None}
        时跳过该半边，不 panic；
  \item \texttt{is\_orientable}：边界半边（\texttt{twin.face = None}）跳过，
        不视为矛盾；
  \item \texttt{fix\_orientations}：对 \texttt{face\_ids().collect::<Vec<\_>>()}
        先收集所有面 ID，避免借用冲突。
\end{itemize}

\section{测试覆盖}

\begin{center}
\renewcommand{\arraystretch}{1.18}
\begin{tabular}{@{}ll@{}}
  \toprule
  \textbf{测试名} & \textbf{验证点} \\
  \midrule
  \texttt{icosphere\_has\_consistent\_normals}    & icosphere(1) 一致且可定向 \\
  \texttt{fix\_orientations\_noop\_on\_icosphere} & 一致网格翻转数 $= 0$ \\
  \bottomrule
\end{tabular}
\end{center}

全模块共 2 个单元测试，覆盖「干净网格」正常路径；
项目整体 \texttt{cargo test} 共 371 个用例。
\textbf{待补充}：不一致网格翻转修复测试、不可定向网格检测测试（Möbius 构造）。

\section{设计权衡}

\begin{itemize}
  \item \textbf{BFS 而非 DFS}：避免栈溢出，对超大连通分量更稳定；
  \item \textbf{仅支持三角面翻转}：本项目所有网格为三角网格，简化实现；
        若未来需要支持多边形，需扩展 \texttt{flip\_face\_orientation}；
  \item \textbf{fix\_orientations 不返回错误}：不可定向网格无法修复，
        但本函数仍尝试修复可修复的部分，返回翻转数；
        调用方可通过后续 \texttt{are\_normals\_consistent} 检查是否完全修复；
  \item \textbf{visited 用 HashSet}：\texttt{FaceId} 实现 \texttt{Hash}，
        查询 $O(1)$。
\end{itemize}

\end{document}
